Let $K$ be a field and let
\[
  \Omega=\{+,-,\times,/\}
\]
be the ordinary binary arithmetic signature.  Division is partial in the
affine chart: a denominator must be nonzero at the moment at which the
operation is performed.

\subsection{Chronological comb conventions}

\begin{definition}[Left- and right-expanded one-hole expressions]
\label{def:p0-left-right-combs}
Fix constants $c_1,\ldots,c_n\in K$ and operation labels
$\omega_1,\ldots,\omega_n\in\Omega$.  Starting from the hole $z$, define
recursively
\begin{equation}\label{eq:p0-comb-recursions}
  L_0[z]=z,
  \qquad
  L_k[z]=\omega_k(L_{k-1}[z],c_k),
\end{equation}
\begin{equation}\label{eq:p0-right-comb-recursions}
  R_0[z]=z,
  \qquad
  R_k[z]=\omega_k(c_k,R_{k-1}[z]).
\end{equation}
The expression $L_n$ is a \emph{left-expanded comb}; $R_n$ is a
\emph{right-expanded comb}.  The index $k$ is the evaluation time: the
innermost operation has index $1$ and is evaluated first.
\end{definition}

Thus, in fully parenthesized form,
\[
  L_n[z]
  =\omega_n(\cdots\omega_2(\omega_1(z,c_1),c_2)\cdots,c_n),
\]
whereas
\[
  R_n[z]
  =\omega_n(c_n,\omega_{n-1}(c_{n-1},\cdots\omega_1(c_1,z)\cdots)).
\]
In the first expression the active subtree is the left child at every internal
vertex.  In the second it is the right child.  Both trees have one internal
spine.  The terminology refers to that explicit slot convention, not to the
visual direction in which a typesetter happens to draw the tree.

\begin{figure}[t]
\centering
\begin{tikzpicture}[
  op/.style={circle,draw=black!65,fill=blue!7,minimum size=7mm,inner sep=0pt},
  leaf/.style={rectangle,draw=black!45,fill=gray!5,minimum height=6mm,
              minimum width=8mm,inner sep=1pt},
  every node/.style={font=\small}
]
  \node[font=\bfseries] at (-3.2,2.65) {left-expanded};
  \node[op] (l3) at (-3.2,2.0) {$\omega_3$};
  \node[op] (l2) at (-4.0,1.15) {$\omega_2$};
  \node[leaf] (lc3) at (-2.35,1.15) {$c_3$};
  \node[op] (l1) at (-4.75,0.3) {$\omega_1$};
  \node[leaf] (lc2) at (-3.25,0.3) {$c_2$};
  \node[leaf] (lz) at (-5.35,-0.55) {$z$};
  \node[leaf] (lc1) at (-4.15,-0.55) {$c_1$};
  \draw (l3)--(l2); \draw (l3)--(lc3);
  \draw (l2)--(l1); \draw (l2)--(lc2);
  \draw (l1)--(lz); \draw (l1)--(lc1);
  \draw[-{Latex[length=2mm]},very thick,blue!60!black]
    (-5.75,-0.9) .. controls (-5.2,-1.25) and (-3.6,-1.25) .. (-2.8,-0.9)
    node[midway,below,font=\scriptsize] {evaluation follows the left spine outward};

  \node[font=\bfseries] at (3.2,2.65) {right-expanded};
  \node[op] (r3) at (3.2,2.0) {$\omega_3$};
  \node[leaf] (rc3) at (2.35,1.15) {$c_3$};
  \node[op] (r2) at (4.0,1.15) {$\omega_2$};
  \node[leaf] (rc2) at (3.25,0.3) {$c_2$};
  \node[op] (r1) at (4.75,0.3) {$\omega_1$};
  \node[leaf] (rc1) at (4.15,-0.55) {$c_1$};
  \node[leaf] (rz) at (5.35,-0.55) {$z$};
  \draw (r3)--(rc3); \draw (r3)--(r2);
  \draw (r2)--(rc2); \draw (r2)--(r1);
  \draw (r1)--(rc1); \draw (r1)--(rz);
  \draw[-{Latex[length=2mm]},very thick,orange!75!black]
    (5.75,-0.9) .. controls (5.2,-1.25) and (3.6,-1.25) .. (2.8,-0.9)
    node[midway,below,font=\scriptsize] {evaluation follows the right spine outward};
\end{tikzpicture}
\caption{The two pure combs under the chronological indexing convention.
Both have one chain of internal dependencies; the active subtree occupies
opposite operand slots.}
\label{fig:p0-two-combs}
\end{figure}

\begin{proposition}[Unique internal order]
\label{prop:p0-combs-unique-order}
Every finite left- or right-expanded comb has a unique legal order of internal
evaluation, namely the chronological order $1,2,\ldots,n$.
\end{proposition}

\begin{proof}
At stage $k$, the only internal child of the $k$th vertex is the subtree formed
at stage $k-1$; the other child is the leaf $c_k$.  Hence the internal vertices
form one dependency chain
\[
  v_1<v_2<\cdots<v_n.
\]
A finite chain has a unique linear extension.
\end{proof}

This proposition is deliberately weaker than the full sequential-tree
classification of Paper~I.  Here the combs are given by explicit grammars;
Paper~I proves intrinsically that every binary expression tree with a unique
internal evaluation order has one such spine after operand slots are recorded.

\subsection{One-hole sections and opposite operations}

For $\omega\in\Omega$ and $c\in K$, write
\[
  C^{(1)}_{\omega,c}[z]=\omega(z,c),
  \qquad
  C^{(2)}_{\omega,c}[z]=\omega(c,z).
\]
Then $L_n$ is a chronological word of pure slot-$(1)$ sections and $R_n$ a
word of pure slot-$(2)$ sections.

Define the opposite binary operation by
\[
  \omega^{\mathrm{op}}(u,v):=\omega(v,u).
\]

\begin{proposition}[Comb opposition]
\label{prop:p0-comb-opposition}
With chronological indexing, a right-expanded comb for operations
$\omega_1,\ldots,\omega_n$ is exactly a left-expanded comb for the opposite
operations:
\[
  R_k[z]
  =\omega_k^{\mathrm{op}}(R_{k-1}[z],c_k).
\]
Consequently, the two comb grammars are isomorphic after replacing every
operation by its opposite.  They agree literally for commutative operations.
\end{proposition}

\begin{proof}
By definition,
\[
  \omega_k(c_k,R_{k-1}[z])
  =\omega_k^{\mathrm{op}}(R_{k-1}[z],c_k).
\]
Apply this identity at every stage.  If $\omega_k$ is commutative, then
$\omega_k^{\mathrm{op}}=\omega_k$.
\end{proof}

The result is combinatorial, not an assertion that the two arithmetic maps are
always equal.  Their elementary one-hole maps are:

\begin{center}
\renewcommand{\arraystretch}{1.35}
\begin{tabular}{lll}
\toprule
Operation & active value in slot $(1)$ & active value in slot $(2)$ \\
\midrule
addition & $z+c$ & $c+z=z+c$ \\
subtraction & $z-c$ & $c-z$ \\
multiplication & $cz$ & $cz$ \\
division & $z/c$ & $c/z$ \\
\bottomrule
\end{tabular}
\end{center}

For $c\ne0$, every first-slot map is affine.  The second-slot subtraction is
also affine, but it has multiplier $-1$.  The second-slot division is not
affine: it has a finite pole at $z=0$ and sends that point to infinity under
projective continuation.

\begin{corollary}[Where the algebraic symmetry breaks]
\label{cor:p0-affine-projective-break}
If the operation alphabet is restricted to $+$ and $\times$, pure left and
pure right expansion have the same one-hole operator algebra.  With
subtraction they remain inside the affine group.  Once second-slot division
$c/z$ is admitted, the right-expanded language contains inversion and leaves
the affine stabilizer of infinity.
\end{corollary}

\begin{proof}
The first two assertions follow from the table.  For $c\ne0$,
\[
  \frac{c}{z}=D_c\circ J(z),
  \qquad D_c(z)=cz,\quad J(z)=\frac1z.
\]
The map $J$ exchanges $0$ and $\infty$, so it does not stabilize infinity.
\end{proof}

\subsection{Point propagation and condition propagation}

The same one-hole word supports two elementary readings.  Given maps
$F_1,\ldots,F_n$ and a point $z_0$, the prefix operators
\[
  G_k=F_k\circ\cdots\circ F_1
\]
push the point forward:
\[
  z_k=G_k(z_0).
\]
This is the \emph{point propagation} used for paths.

Alternatively, given a family of output level sets $\{\Sigma_t\}$, the same
prefix operator pulls each level set back:
\[
  \Sigma_{k,t}=G_k^{-1}(\Sigma_t).
\]
This is \emph{condition propagation}.  It is contravariant: an outer condition
is transported toward the input by inverse image.  Sections~\ref{sec:p0-paths}
and~\ref{sec:p0-ripples} show that the first reading yields a zigzag path,
while the second yields a ripple pencil.
